\section{Data and Scoring}

\subsection{The ROUTE Corpus}

The ROUTE corpus consists of 227 interstate corridors constructed from the FHWA TIGER/Line 2023 Primary Roads national shapefile \citep{Census_TIGER2023}. Each corridor is defined as all road segments sharing a common route identifier (e.g., \texttt{I80}, \texttt{I95}) in the TIGER data. The corpus covers 48,792 miles of interstate highway across 5,599 individual segments.

For each corridor, we compute or join the following attributes from federal sources:

\begin{itemize}
\item \textbf{AADT and truck fraction}: from FHWA Highway Performance Monitoring System (HPMS) 2018 data \citep{FHWA_HPMS2018}, accessed via the FHWA geo.dot.gov ArcGIS REST service. We fetched 160,969 records for interstate-class roads (functional system code 1, \texttt{f\_system = 1}) across 28 states that returned data. Truck fraction is computed as \texttt{aadt\_combination / aadt} (combination truck count divided by total AADT). For each route, we aggregate across states using the median AADT (robust to outlier state submissions).

\item \textbf{Lane count and IRI}: also from HPMS 2018. Lane count determines theoretical daily capacity ($\text{lanes} \times 1{,}900 \text{ pcph} \times 24\text{h}$). Theoretical daily capacity is computed as: $C = n_{\text{lanes}} \times 1{,}900 \text{ pcphpl} \times 24\text{h} \times K^{-1}$, where $K = 0.09$ is the peak-hour factor, yielding approximately 37{,}000 vehicles per day per lane-pair per direction \citep{HCM7}. IRI (International Roughness Index) is the primary pavement quality measure.

\item \textbf{Bridge condition}: from FHWA National Bridge Inventory 2023 \citep{FHWA_NBI2023}, joined by coordinate proximity (tolerance 0.002$^\circ$, approximately 170m, with route-name similarity check).

\item \textbf{Freight commodity flows}: from FAF5 v5.6 \citep{FAF5_2022}. Flows are attributed to corridors using a zone-traversal method: we sum all truck-mode flows between FAF origin-destination pairs whose centroids fall within the corridor's 50-mile buffer. This is an approximation; routing-based attribution is deferred to future work.

\item \textbf{Network centrality}: computed using the Brandes algorithm \citep{Brandes2001} on the national directed highway graph (12,421 edges, 12,011 nodes, including US highway segments used as supplementary context). Betweenness centrality scores are normalized to [0, 1] across all edges; corridor-level centrality is the mean of edge-level scores.
\end{itemize}

\paragraph{B2 Reliability Caveat.}
The B2 centrality scores are computed via a simplified Brandes algorithm on a
partial directed graph derived from the TIGER/Line 2023 Primary Roads dataset
\citep{Census_TIGER2023}. Two sources of approximation should be noted. First, the
graph is partial: TIGER/Line captures primary roads but may miss secondary
connections that alter centrality rankings on less-connected corridors. Second,
the Brandes implementation uses a simplified predecessor tracking that may
affect absolute centrality values. However, the claim of the paper is not that
absolute B2 values are precise but that \textbf{relative rankings} among the 8
highest-centrality corridors are directionally correct. Three independent
validations support this: (a) the 8 T1 corridors are also the 8 highest-degree
nodes in the national graph (I-5, I-10, I-80, I-90, I-95 each have 5+ T1/T2
intersections; I-110 has 2), confirming topology-based identification independent
of Brandes; (b) STRAHNET designation aligns at 100\%; (c) ATRI bottleneck density
ranks the same 8 corridors highest at $\rho = 0.72$ \citep{ATRI_costs2024}.

\subsection{The 12-Dimension Pool}

We score each corridor against 12 candidate dimensions across four bands. Table~\ref{tab:dimensions} provides the full dimension definitions and scoring anchors.

\begin{table}[h!]
\centering
\caption{12-Dimension Scoring Pool (ROUTE v1.0 Rubric)}
\label{tab:dimensions}
\small
\begin{tabular}{llp{5.5cm}}
\toprule
\textbf{Dim} & \textbf{Name} & \textbf{Primary data source} \\
\midrule
A1 & Throughput Gap & HPMS AADT vs. lane-based capacity \\
A2 & Freight Intensity & FAF5 commodity value, HPMS truck fraction \\
A3 & Speed Reliability & FHWA Freight Performance Measures (PTI); IRI proxy when unavailable \\
\midrule
B1 & Redundancy & Graph shortest-path; detour penalty miles \\
B2 & Network Centrality & Brandes betweenness on HighwayGraph \\
B3 & Port/Border Access & BTS port rankings; FHWA border crossings \\
\midrule
C1 & Population Reach & Census ACS 2022 (50-mile buffer) \\
C2 & Rural Connectivity & USDA ERS rural codes; interchange gap analysis \\
C3 & Economic Opportunity & BEA GDP per capita (buffer relative to national) \\
\midrule
D1 & Climate Resilience & FEMA NFHL flood zones (miles in SFHA) \\
D2 & Multimodal Integration & AAR intermodal hubs; DOE AFDC EV charger density \\
D3 & Infrastructure Vintage & NBI condition ratings; HPMS IRI; bridge age \\
\bottomrule
\end{tabular}
\end{table}

All dimensions are scored 0--10 using a three-point linear interpolation on published scoring anchors (see Appendix~A for full anchor table). Anchors are runtime-configurable rather than compiled constants, following the calibration-first design of the ROUTE system: anchors will be rescaled after A.2's calibration pass once 20+ corridors have been scored under stable data conditions.

\subsection{Data Availability and Known Limitations}

\textbf{HPMS coverage}: 28 of 50 states returned usable data from the geo.dot.gov service during our fetch window. States without HPMS coverage (including TX, TN, VA, WY, and others) receive estimated A1/A2/A3 scores marked with a dagger (†). This introduces systematic bias: some major corridors (I-10 in Texas, I-81 in Virginia) are scored from partial data.

\textbf{FAF5 attribution}: our zone-traversal method over-counts flows that use multiple parallel corridors within a FAF zone and under-counts through-flows that do not terminate within the buffer. All A2 scores are marked estimated.

\textbf{Speed reliability}: FHWA Freight Performance Measures (PTI data) were not available at time of analysis; A3 scores fall back to an IRI-based proxy, which captures pavement roughness but not congestion-driven unreliability. This creates a known bias toward higher A3 scores on rough-pavement corridors.

Despite these limitations, the 227-corridor scoring represents the most comprehensive publicly-available data-driven assessment of the US interstate network to date.
